\section{Introducción}
Todos los dispositivos electrónicos activos requieren una fuente de corriente
directa (CD) constante que provenga de una batería o una fuente de
alimentación. La \textbf{fuente de alimentación de CD} convierte el voltaje de
corriente alterna (CA) estándar de $220[\text{V}]$ a $50[\text{Hz}]$ disponible
en las tomas de corriente de pared en un voltaje de CD constante.

En la \textbf{figura~\ref{diagrama}} se muestra un diagrama de bloques básico de
una fuente de alimentación completa.

\begin{figure}[!h]
\centering
\includegraphics[scale=1.50]{diagramas/00.diagrama.eps}
\caption{Fuente de alimentación completa.}
\label{diagrama}
\end{figure}

En general, el voltaje de linea de entrada de CA se reduce a un voltaje de CA
más bajo con un \textbf{transformador}. Este cambia voltajes de CA con base en
la relación de vueltas entre el primario y el secundario. Si éste tiene más
vueltas que el primario, el voltaje de salida a través del secundario será más
alto y la corriente será más pequeña. Si el secundario tiene menos vueltas que
el primario, el voltaje de salida a través del secundario será más bajo y la
corriente será más alta.

El \textbf{rectificador} puede ser de media onda o de onda completa, este
convierte el voltaje de entrada de CA en un voltaje de CD pulsante.

El \textbf{filtro} elimina los rizos de voltaje en el rectificador y produce un
voltaje de CD relativamente uniforme.

El \textbf{regulador} es un circuito que mantiene un voltaje de CD constante
frente a las variaciones de voltaje de linea de entrada o de la carga. Los
reguladores varían desde un dispositivo de un solo semiconductor hasta circuitos
integrados mas complejos.

La \textbf{carga} es un circuito o dispositivo conectado a la salida de la
fuente de alimentación y opera con el voltaje y la corriente de la fuente de
alimentación \cite{Floyd}.

